\section{Auction-Based Multi-Policy Framework for Multi-Objective RL}

\textbf{Preliminaries: multi-objective RL.}
We use multi-objective Markov decision processes (\momdp) to model uncertain environments and multiple objectives. 
An \momdp with \(m \in \N_{>0}\) objectives is specified by a tuple \(\mathcal{M} = (S, A, T, R, \gamma, \mu_0)\) where \(S\) is the set of states, \(A\) is the set of actions, \(T : S \times A \times S \to \R\) is the transition function so that \(s' \mapsto T(s, a, s')\) forms a probability distribution over \(S\) for all \((s,a) \in S \times A\), \(\mathbf{R}=\{R_i \colon S \times A \times S \to \R\}_{i\in [1;m]}\) is the set of reward functions corresponding to the \(m\) objectives, $\gamma\in (0,1)$ is the discount factor, and \(\mu_0 \in \dist(S)\) is the initial state distribution. 
A \textit{policy} in an \momdp~\(\mathcal{M}\) is a function \(\pi : (S \times A)^* \times S \to \dist(A)\) that maps a history of state-action pairs and the current state to a distribution over actions.
The policy $\pi$ gives rise to a random sequence of states and actions $s^0,a^0,s^1,a^1,\ldots$, that produces a sequence of rewards $r^0_ir^1_i\ldots$ for each objective $i\in [1;m]$, where $r^t_i = R_i(s^t,a^t,a^{t+1})$ for every $t\in \mathbb{N}$.
The \textit{value} of $\pi$, denoted $V^\pi$, is a vector in $\R^m$ such that $V^\pi_i = \E_{s_t,a_t}[ \sum_{t=0}^\infty \gamma^t r^t_i]$.
In the multi-objective setting, a commonly studied problem is to find a policy in the pareto-frontier of the value function which balances the trade-offs between the objectives~\citep{van2014multi}.

\textbf{The auction-based multi-policy framework.}
We consider the compositional approach to policy synthesis for \momdps, where we will design a selfish, \emph{local} policy maximizing each individual value component, towards the fulfillment of some required global coordination requirements.
The main crux is in the composition process, where each local policy may propose a different action, but the composition must decide one of the actions that will be actually executed.
Importantly, the composition must be implementable in a distributed manner, meaning we will \emph{not} use any global policy that would pick an action by analyzing all local policies and their reward functions.

We present a novel \textit{auction}-based RL framework for compositional policy synthesis for \momdps.
In our framework, not only do the local policies emit actions, but they also \emph{bid} for the privilege of executing their actions for a given number of time steps $\tau\in \mathbb{Z}_{>0}$ in future.
The bids are all non-negative integers bounded by a given constant $\beta\in \mathbb{Z}_{\geq 0}$, and the highest bidder's actions get executed for the following $\tau$ consecutive steps, with bidding ties being resolved uniformly at random.
To discourage policies from always bidding high, we penalize the policies proportional to their bids scaled with a constant $\rho\in (0,1)$.
We consider two intuitive alternatives for the penalty model, namely \poorman and \allpay:
in the \poorman setting, only the executed policy (one of the highest bidders) pays the penalty, whereas in the \allpay setting, all policies pay penalties.
Through this mechanism, policies select their bids to reflect how critical their action is in the current state, and the bidding penalty incentivizes truthful urgency.
This way, we obtain a purely decentralized scheme to coordinate local policies in a given \momdp.

The parameters $\tau$, $\beta$, and $\rho$ denote the \emph{bidding interval}, \emph{maximum bid}, and \emph{bid penalty coefficient}, respectively. Ablation studies in Section~\ref{sec:experimental-results} highlight their importance. If $\tau$ is too small, policies switch too frequently, risking suboptimal outcomes. If $\beta$ is too small, policies cannot express preference strength, leading to near-random control. If $\rho$ is too large, policies tend to abstain from bidding.

\textbf{Local policies via multiplayer Markov games.}
We formalize our framework as an $m$-player general-sum Markov game~\citep{littman1994markov}, where each player optimizes one of the $m$ objectives of the original \momdp. The game encodes the bidding mechanism in its transitions and rewards, and we therefore call it a \textit{Markov bidding game}. Each state records the current \momdp state, the policy selected in the last bidding round, and the remaining time until the next bidding. At each state, players independently choose an action and a bid; transitions determine the winning policy based on the bids, while rewards capture the corresponding bidding penalties. 


\begin{definition}[Markov bidding games]\label{def:markov game}
    Let \(\mathcal{M} = (S, A, T, R, \gamma, \mu_0)\) be an \momdp with \(m\) objectives, $\tau$ be the bidding interval, $\beta$ be the maximum bid, $\rho$ be the bid penalty coefficient, and $\Delta\in \{ \poorman,\allpay\}$ be the used penalty model.
    We define the \(m\)-player Markov bidding game as a tuple \(\mathcal{G}^{\mathcal{M}}_{\tau,\beta,\rho} = (\hat S, \hat{A}, \hat{T}, \hat{R}, \gamma, \hat \mu_0)\) where
    \begin{itemize}
        \item \(\hat S = S \times [1;m] \times [0;\tau - 1]\) is the state space. An element of \(\hat S\) is denoted \((s, k, t)\).
        \item \(\hat{A} = (A \times \Z_{\geq 0})^m\) represents the joint extended action space of the \(m\) players, where each player independently selects an action from \(A\) and a nonnegative bid from $\Z_{\geq 0}$. We denote a set of actions chosen by the \(m\) players with \((\vec a, \vec b) = (a_i, b_i)_{i \in [1;m]} \in \hat{A}\).
        \item \(\hat{T} : \hat S \times \hat{A} \times \hat S \to \R\) is the new transition function defined as,
        \begin{equation*}
            \hat T((s, k, t), (\vec a, \vec b), (s', k', t')) \coloneqq 
            \begin{cases}
                T(s, a_{k'}, s') & t >0\land t' = t-1 \land k'=k,  \\
                \frac{1}{\abs{\vec{b}_{\max}}} T(s, a_{k'}, s') & t = 0 \wedge t' = \tau-1 \wedge k' \in \vec{b}_{\max}, \\
                0 & \text{otherwise},
            \end{cases}
        \end{equation*}
        where \(\vec b_{\max} \coloneqq \{i \mid b_i = \max \{b_1, \dots, b_m\}\}\) is the set of agents with the highest bid.
        \item \(\{\hat{R}_i : \hat S \times \hat{A} \times \hat S \to \R\}_{i \in [1; m]}\) is the set of reward functions for the \(m\) agents with
        \begin{equation*}
            \hat{R}_i((s, k, t), (\vec a, \vec b), (s', k', t')) \coloneqq 
            \begin{cases}
                R_i(s, a_{k'}, s') - \rho \cdot b_{i} & t = 0 \land \left( i = k' \lor \Delta = \allpay \right),\\
                R_i(s, a_{k'}, s') & \text{otherwise}.
            \end{cases}
        \end{equation*}
        \item $\hat \mu_0$ is the initial state distribution over \(\hat S\) defined as,
        \begin{equation*}
        	\hat \mu_0(s,k,t) \coloneqq
        		\begin{cases}
        			\mu_0(s)	&	t=0 \land k=i, \\
        			0			&	\text{otherwise},
        		\end{cases}
        \end{equation*}
        where $i\in [1;m]$ is an arbitrary agent.
    \end{itemize}
\end{definition}

In the Markov game $\mathcal{G}^{\mathcal{M}}_{\tau,\beta,\rho}$, each player $i$ aims to compute its own selfish local policy of the form $\pi_i\colon (\hat{S}\times\hat{A})^*\times \hat{S}\mapsto (a_i,b_i)\in A\times \mathbb{Z}_{\geq 0}$, such that the discounted sum of its own reward is maximized, given that the other players are selfishly fulfilling their own objectives.

\textbf{Learning policies in Markov bidding games.}
Since the transitions of the original \momdp are unknown, the transitions of the constructed Markov game in Definition~\ref{def:markov game} will also be unknown.
Our goal is to learn the local policy of each player using RL.
\citet{yu2022mappo} demonstrate that proximal policy optimization (PPO; \cite{schulman2017ppo}) can be effective to train multiple policies concurrently for multi-agent settings. We follow their recipe and train all the local policies with PPO. 
Similar to \citet{yu2022mappo}, since the objectives in our environments share the same characteristics, we use the same actor and critic network for all the policies. 
At each iteration of PPO, we collect all the rollouts from each policy's perspective and sample minibatches from this pool to perform the PPO updates.

Additionally, to cope with varying number of objectives and varying number of players in the resulting Markov game, we employ the \textit{attention pooling} architecture to compress any number of objectives into a fixed-sized embedding. 
This design enables the agent to operate with an arbitrary number of policies (i.e., objectives) at deployment time. 
Each objective information vector is first processed by a fully connected neural network to produce a target embedding. 
These embeddings are then aggregated through a weighted sum, where the weights are computed by taking the dot product between each embedding and a learned query vector. 
This attention mechanism produces a fixed-dimensional representation that summarizes the set of objectives while allowing the model to adapt to varying objective counts.


% \section{A Multi-Agent Bidding Approach for Multi-Objective RL}
% \begin{definition}[\momdp]
%     A multi-objective Markov decision process (\momdp) with \(m \in \Z_{>0}\) objectives is specified by a tuple \(\mathcal{M} = (S, A, T, R, \mu_0)\), where
%     \begin{itemize}
%         \item \(S\) is the set of states,
%         \item \(A\) is the set of actions,
%         \item \(T : S \times A \to \dist(S)\) is the transition function mapping a state-action pair to a distribution over states,
%         \item \(R : S \times A \times S \to \R^m\) is the reward function with each output component corresponding to the different objectives, and
%         \item \(\mu_0 \in \dist(S)\) is the initial state distribution.
%     \end{itemize}
% \end{definition}

% A \textit{policy} in an \momdp~\(\mathcal{M}\) is a function \(\pi : (S \times A)^* \times S \to \dist(A)\) that maps a history of state-action pairs and the current state to a distribution over actions. 

% \begin{definition}[\mabmdp]
%     Let \(\mathcal{M} = (S, A, T, R, \mu_0)\) be an \momdp~with \(m\) objectives and let \(b \in \Z_{>0}\) be the bid upper bound. Also, define \(M = \{1, \dots, m\}\) be indices of the \(m\) agents corresponding to the \(m\) objectives along with \(\bot\) representing a null agent. Lastly, let \(B = \{0, \dots, b\}\) be the range of bids and \(\rho > 0\) be the bid penalty factor.
%     We define the multi-agent bidding Markov decision process (\mabmdp) as a tuple \(\mathcal{B}_{\mathcal{M}} = (\hat S, \hat{A}, \hat{T}, P, \hat{R}, \hat \mu_0)\) where
%     \begin{itemize}
%         \item \(\hat S = M \times S\) is the new state space augmented with the index of the agent that won the previous round of bidding,
%         \item \(\hat{A} = A^m \times B^m\) represents the action space of the \(m\) agents in which each agent selects an action from \(A\) and a bid from \(B\),
%         \item \(\hat{T} : \hat S \times \hat{A} \to \dist(\hat S)\) is the new transition function defined as,
%         \begin{equation*}
%             \hat{T}((\_, s), (\vec a, \vec b)) \coloneqq \frac{1}{\abs{B_{\max}}}\sum_{i \in B_{\max}} (T(s, a_i), i)
%         \end{equation*}
%         where \(B_{\max} \coloneqq \{i \mid b_i = \max \{b_1, \dots, b_m\}\}\) is the set of agents with maximal bids. The tuple \((T(s, a_i), i)\) represents the distribution over \(\hat S\) induced by the original transition function \(T\) such that the second component is fixed, and the weighted sum represents taking the weighted sums of the distributions over \(\hat S\).
%         \item \(P : \hat A \times M \to \R^m\) is the bidding penalty for the \(m\) agents and the second component is the index of the agent that won the bidding.
%         \item \(\hat{R} : \hat S \times \hat{A} \times \hat S \to \R^m\) is the reward function for the \(m\) agents with
%         \begin{equation*}
%             \hat{R}_k((\_, s_0), (\vec a, \vec b), (i, s)) \coloneqq R_k(s_0, a_i, s) - P_k((\vec a, \vec b), i)
%         \end{equation*}
%         where \(i \in M\) is the index of the agent that won the bid and chose the action.
%         \item \(\hat \mu_0 \coloneqq (\mu_0, 1)\) is the initial state distribution over \(\hat S\) induced by \(\mu_0\) and the second component is fixed to be \(1\) without loss of generality.
%     \end{itemize}
% \end{definition}

% Given an \mabmdp~\(\mathcal{B}_{\mathcal{M}}\), a \textit{policy} for each agent indexed by \(i \in \{1, \dots, m\}\) takes a similar form: \(\pi_i : (\hat S \times \hat A)^* \times \hat S \to \hat A\). Intuitively, a state \((i, s) \in \hat S\) encodes the agent that won the bidding and chose the action to reach \(s\) in the previous step. At each step, each of the agents choose an action and a bid, and an action amongst the set of highest bidders is chosen uniformly at random. The reward function includes a penalty term that captures the desired bidding mechanism.
